\section{A genuine family on which the variational exclusion is nondecisive}
\label{sec:s5-endpoint-scope}

The exact all-prime inequality in MA-RB-008 narrows an actual
comparison obstruction. Its numerical exclusion criterion
cannot hold uniformly on Calegari's family. We prove this
using genuine fields satisfying both specified local
hypotheses. No zero of their \(L\)-functions is constructed.

\paragraph{MA-RB-009 (scope of the universal numerical criterion).}
Retain the final discriminant-aware functions
\(\mathcal H_\Delta,\mathcal F_{\Delta,k},\mathcal J_{\Delta,k}\)
of MA-RB-008, and put
\[
 C_0=60(1+\log(2\pi))+5\log5+\frac{40}{7},
 \qquad C_\Delta=\frac12\log\Delta-C_0 .
 \tag{MA-RB.97}
\]
For every actual field in the specified family and every
integer \(k\ge1\), if \(C_\Delta\ge18k\), then
\[
 \frac14-\frac{9k}{C_\Delta}
       \le\mathcal J_{\Delta,k}<\frac14.
 \qquad
 C_\Delta\ge40k\ \Longrightarrow\
                      \mathcal J_{\Delta,k}\ge\frac1{40}>0 .
 \tag{MA-RB.98}
\]

There are genuine fields in this family with arbitrarily
large normal-closure discriminant. An explicit construction
is as follows. For each rational prime \(q>5\), choose
the least positive solution \(a_q\) to the two congruences
\[
 qa_q\equiv2\pmod3,\qquad qa_q\equiv1\pmod5,
 \qquad
 f_q(X)=X^5+qa_qX+15q .
 \tag{MA-RB.99}
\]
Let \(K_q=\mathbb Q[X]/(f_q)\) and let \(L_q\) be its
normal closure. Then
\[
 \begin{gathered}
 \operatorname{Gal}(L_q/\mathbb Q)=S_5,\qquad
 [L_q:\mathbb Q]=120,\qquad
 (r_1(L_q),r_2(L_q))=(0,60),\\
 \text{complex conjugation has type }(12)(34),\\
 5\text{ is unramified and }
       \operatorname{Frob}_5\text{ has type }(12)(34),\qquad
 q\mid\Delta_{L_q}.
 \end{gathered}
 \tag{MA-RB.100}
\]
Consequently, for every fixed integer \(k\ge1\),
\[
                \mathcal J_{\Delta_{L_q},k}
                        \longrightarrow\frac14
                   \qquad(q\longrightarrow\infty
                         \text{ through primes}).
 \tag{MA-RB.101}
\]
In particular the criterion
\(\mathcal J_{\Delta,1}\le0\) does not exclude comparison
failure uniformly for this actual family.
This is a limitation of that criterion, not a construction
of comparison failure, a shared real zero, or a
\(\beta\)-excess example. It does not reduce the remaining
shared-zero gate.

\begin{proof}
\emph{Construction, irreducibility, and real signature.}
Because \(q\) is prime to \(15\), the Chinese remainder
theorem supplies \(a_q\), with \(1\le a_q\le15\).
In particular \(qa_q>0\).
The polynomial \(f_q\) is Eisenstein at \(q\):
every nonleading coefficient is divisible by \(q\),
and its constant coefficient has \(q\)-valuation one.
Here the irreducibility argument is elementary.
A factorization into two monic integer polynomials
of positive degree, which a rational factorization of
a monic polynomial would supply, reduces modulo \(q\)
to a factorization of \(X^5\). Both constant terms
would be divisible by \(q\), contrary to
\(q^2\nmid15q\).
Thus \(f_q\) is irreducible over \(\mathbb Q\), and
its Galois group acts transitively on its five roots.

On the real line,
\[
                         f_q'(x)=5x^4+qa_q>0.
\]
The polynomial tends to opposite infinities at the
two ends of that line. Hence it has exactly one
real root. Characteristic zero and irreducibility
give separability; the other four roots therefore
form two nonreal conjugate pairs. Complex conjugation
on the five roots has cycle type \(2B=(12)(34)\).

\emph{The two finite Frobenius calculations.}
The defining congruences give the exact factorizations
\[
 \begin{aligned}
 f_q(X)\bmod3
   &=X^5-X=X(X-1)(X+1)(X^2+1),\\
 f_q(X)\bmod5
   &=X^5+X=X(X^2+2)(X^2-2).
 \end{aligned}
 \tag{MA-RB.102}
\]
Over \(\mathbb F_3\), the first three linear factors
are distinct, and \(X^2+1\) is irreducible, since
\(-1\) is not a square. All four factors are distinct.
Over \(\mathbb F_5\), neither \(2\) nor \(-2=3\)
is a square, since the nonzero squares are \(1,4\).
The two quadratic factors and \(X\) are distinct.
Both reductions are therefore squarefree.

We spell out the local interpretation being used.
The roots of either reduction belong to a finite
residue-field extension and lift uniquely, by the
simple-root Hensel property, in the corresponding
unramified local extension. All characteristic-zero
roots of \(f_q\) lie there: there are five distinct
lifts and the polynomial has degree five.
Thus the local splitting extension is unramified
at \(3\) and at \(5\). The arithmetic residue
Frobenius permutes the reduced roots in cycles whose
lengths are their irreducible factor degrees.
Uniqueness of the lifts gives the same permutation
on the characteristic-zero roots.
This proves that \(3\) supplies a transposition,
while \(5\) supplies the specified \(2B\) element.
It also proves the required unramified assertion
in the normal closure, rather than merely in
the degree-five subfield.

\emph{Why the transitive group is all of \(S_5\).}
Let \(G\) be that transitive group, and let
\(\tau\in G\) be the transposition just obtained.
On the five roots, form a graph whose edges are
the two-element supports of all \(G\)-conjugates
of \(\tau\). The graph has an edge and is
\(G\)-invariant. Transitivity implies that all
its connected components have the same size.
That size divides the prime number \(5\), and
it is not one because the graph has an edge.
The graph is therefore connected.
The transpositions associated with the edges
of a connected graph generate its full symmetric
group: along a path, conjugating an edge transposition
by the preceding edge transpositions produces the
transposition of the first and last vertices.
All those edge transpositions lie in \(G\).
It follows that \(G=S_5\).

In particular the normal closure has degree \(120\).
Its complex conjugation is nontrivial. A Galois
number field with a real embedding would, after
that embedding, have every conjugate in \(\mathbb R\)
and hence trivial complex conjugation.
Thus \(L_q\) is totally imaginary, giving signature
\((0,60)\).

\emph{Ramification at the varying prime and unbounded discriminants.}
Extend the \(q\)-adic valuation to \(L_q\), normalized
by \(v(q)=1\), and let \(\theta\) be any root of \(f_q\).
If \(v(\theta)<0\), then \(\theta^5\) has valuation
strictly smaller than both other terms of the
polynomial equation, which is impossible for a sum
equal to zero. The same contradiction holds if
\(v(\theta)=0\).
If \(v(\theta)>0\), the linear term has valuation
at least \(1+v(\theta)>1\), while the constant
term has valuation one. Thus the minimum valuation
can occur twice only when
\[
                   5v(\theta)=1.
 \qquad
 v(\theta)=\frac15,\qquad
                    5\mid e_q(L_q/\mathbb Q).
 \tag{MA-RB.103}
\]
The last conclusion follows from the value group
\(e_q^{-1}\mathbb Z\) of this extended valuation.
Therefore \(q\) is ramified and \(q\mid\Delta_{L_q}\).
Here the standard discriminant criterion has its
literal finite-algebra meaning: ramification gives
a nonzero nilpotent in \(\mathcal O_{L_q}/q\mathcal O_{L_q}\);
multiplication by its product with any element is
nilpotent and has trace zero. The reduced trace pairing
is degenerate, so its discriminant vanishes modulo \(q\).
Consequently \(\Delta_{L_q}\ge q\), and these
discriminants tend to infinity with \(q\).
There are also infinitely many distinct normal
closures: finitely many number fields have only
finitely many ramified primes in total.
This completes all assertions in (MA-RB.100).

\emph{Uniform control of the retained Euler and gamma terms.}
The coefficients in the positive-prime form of
\(\mathcal B_\Delta\) are at most \(120\).
For all \(\sigma>1\),
\[
 \begin{split}
 0<\mathcal B_\Delta(\sigma)
 &\le120\sum_{n=2}^{\infty}\frac{\log n}{n^6-1}\\
 &\le120\,\frac{64}{63}\sum_{n=2}^{\infty}n^{-5}\\
 &\le120\,\frac{64}{63}
             \left(\frac1{32}+\int_2^\infty x^{-5}\,dx\right)
   =\frac{40}{7}.
 \end{split}
 \tag{MA-RB.104}
\]
The second line uses \(\log n\le n\) and
\(n^6-1\ge(63/64)n^6\).
This bound is uniform in the discriminant and its
ramified prime support. Also the exact known
prime-five term satisfies
\[
          E_5(\sigma)=\frac{120\log5}{5^{2\sigma}-1}
                                    \le5\log5.
\]
The gamma recurrence
\(\psi(\sigma+1)=\psi(\sigma)+1/\sigma\),
and the already proved Jensen lower bound at
\(\sigma+1\), give
\[
 \psi(\sigma)
     \ge\log\sigma-\frac1\sigma\ge-1
                                    \qquad(\sigma>1).
 \tag{MA-RB.105}
\]
Thus the exact denominator in MA-RB-008 obeys
\[
 \begin{aligned}
 \mathcal H_\Delta(\sigma)&\ge C_\Delta,\\
 \mathcal F_{\Delta,k}(\sigma)
   &=x^2-\frac{18kx}{\mathcal H_\Delta(\sigma)}
      \ge x^2-\frac{18kx}{C_\Delta},
          \qquad x=\sigma-\tfrac12>\tfrac12,
 \end{aligned}
 \tag{MA-RB.106}
\]
where the second line is used when \(C_\Delta>0\).
The first line drops only the positive rational
terms \(1/\sigma+1/(\sigma-1)\) and uses
(MA-RB.104)--(MA-RB.105).
If \(C_\Delta\ge18k\), the polynomial on the right
of the second line is increasing for \(x\ge1/2\):
its derivative is \(2x-18k/C_\Delta\ge0\).
Its value at \(x=1/2\) is \(1/4-9k/C_\Delta\).
Taking the infimum and using the already proved
strict upper bound \(\mathcal J_{\Delta,k}<1/4\)
gives (MA-RB.98). At \(C_\Delta\ge40k\) its lower
bound is at least \(1/4-9/40=1/40\).

Finally \(\Delta_{L_q}\to\infty\) implies
\(C_{\Delta_{L_q}}\to\infty\).
For each fixed \(k\), the two bounds in
(MA-RB.98) squeeze its variational minimum
to \(1/4\). This proves (MA-RB.101).
\end{proof}

\paragraph{Actual arithmetic content and limitation.}
The polynomial construction imposes the complex-conjugation
and prime-five hypotheses of Calegari's Theorem~1.1,
assumptions (1)--(2), in the supplied
\texttt{S5ArtinWeb.tex}, lines 123--136.
Calegari's theorem is not being cited as the construction
of this family; that construction and its Galois proof
are given above.
The residue-Frobenius convention is the same arithmetic
convention as (MA-RB.91) and the local regular Artin
factor in MA-RB-008.
The finite-cover placement is the established SGA~1,
Expos\'e~V, \S7 correspondence between finite \'etale
covers and finite continuous fundamental-group sets,
with the Artin-motive interpretation of Deligne (1979),
Definition~6.1 and \S\S6.2--6.3.
The elementary local argument above uses the standard
simple-root Hensel property and the normalized local
valuation; no new theorem locator is assigned to those
classical field-theory facts.

The uniform bounds use the exact Euler expression of
MA-RB-008 and the gamma recurrence and Jensen argument
already proved in MA-RB-007. They require neither
critical-line location, \(\mathsf A\), \(M\), nor
\((\mathcal R_B)\).
The genuine fields demonstrate that one cannot force
the displayed numerical minimum to be nonpositive
for every member merely from these family hypotheses.
They provide no example of the hypothetical comparison
pole or of either shared vanishing factor.
The remaining problem is still an actual arithmetic
shared-zero exclusion on the shorter interval of
MA-RB-008. No new reduction of that interval or
exclusion of its zeros is claimed here.
